%\documentclass{article}
%\usepackage[utf8]{inputenc}
%\begin{document}
%(Type your content here.)
%\end{document}
\documentclass[fleqn,12pt]{uitmkelantan}
%--------------------
\usepackage{amsmath} %utk masuk begin{align}
\usepackage{amsxtra}%supaya \eqref
\usepackage{amsopn}%supaya dapat declare benda baru \citep
\usepackage{amsbsy}
\usepackage{amstext}
\usepackage{amsfonts}
\usepackage{amssymb}
\usepackage{amsthm}
\usepackage{shadow}% kotak shadow
\usepackage{maple2e}% untuk membolehkan kerja maple diekspot 
\usepackage{curves}
\usepackage{natbib}% utk cite (autor,year) guna \citep{}. utk autor(year) guna \cite{} or \citet{}
\usepackage{color}%masukan color
\usepackage{setspace}
%\citestyle{apalike}
%----------------------------
\def\emptyline{\vspace{12pt}}
\DefineParaStyle{Maple Output}
\DefineParaStyle{Maple Plot}
\DefineCharStyle{2D Math}
\DefineCharStyle{2D Output}
%------------------------
\usepackage[utf8]{inputenc}
\usepackage[T1]{fontenc}
\usepackage{times}
\usepackage{lipsum}
\usepackage{marginnote}
\usepackage[top=2.5cm, bottom=2.5cm, outer=2.5cm, inner=3.8cm, heightrounded, marginparwidth=2cm, marginparsep=2cm]
{geometry}
\renewcommand{\ttdefault}{lmtt}
\usepackage{graphicx}
\usepackage{listings}
\lstset{basicstyle=\ttfamily,columns=fullflexible,language=[LaTeX]TeX,frame=single,breaklines=true}
\usepackage{etoolbox}
\usepackage{color}%masukan color
 % ini utk adjust margin
 \oddsidemargin 1.6cm \evensidemargin 0in
%\addtolength{\hoffset}{3.0cm}
\marginparwidth 15pt \marginparsep 7pt
\topmargin 1.0cm %asal1.5
\headsep 0.7cm %asal 1.1
\addtolength{\voffset}{-2.6cm}%adjust 1.25m page no. dari tepi atas
\textheight 20.0cm \textwidth 15.0cm %asal 25.2
%\voffset=0.25cm
\footskip=2cm
\headheight=13pt
 \rightmargin= 3.3cm
 \leftmargin =4.1cm
\textheight=670 pt
 
\AtBeginEnvironment{lstlisting}{\SingleSpacing}

%-------------------------------
\authora{\textrm{Hinnah Aqilah bt Hamzah}}
\studentida{2013276424}
\groupa{}
%-------------------------------
\authorb{Diyana Amalina bt Fadzil}
\studentidb{2013819248}
\groupb{}
%-------------------------------
\authorc{Nor Eleena Binti Hamdan}
\studentidc{2013697362}
\groupc{}
%------------------------------


\supervisor{Prof Madya Dr. Jusoh Yacob}
\code{}
\title{Fresh Graduates Suffer From Underpaid Salary }
\submissionyear{2016}
\degree{Statistics}

% load acronym definitions from separate file
\loadglsentries{sym-acronyms}
%\setlength{\glsdescwidth}{.7\textwidth}
\makeglossaries
%---------------------------
%-----------------------------
\newlength{\defbaselineskip}% kena tambah ini utk spacing
\setlength{\defbaselineskip}{\baselineskip}%sama
\newcommand{\setlinespacing}[1]%
           {\setlength{\baselineskip}{#1 \defbaselineskip}}
           %----------------------------------------------------
\begin{document}
\frontmatter

\makecoverandtitlepage
\clearpage
\thispagestyle{empty}
\chapter{\textrm{DECLARATION}}
\vspace{20mm}
\begin{center}
\textbf{\textrm{APPROVED BY:}}\\
\vspace{5mm}
.................................\\
\textrm{(PROF MADYA DR. JUSOH YACOB)\\
Supervisor
Faculty of Computer and Mathematical Sciences}
\end{center}




\textrm{
\chapter{Acknowledgments} 
Praise be to Allah SWT the Almighty and the All Merciful, who has given us , the powerless creature, His guidance so that we can smoothly finish our final year project entitle "Fresh Graduates Suffer From Underpaid Salary" as the requirement for the degree in Science Statistics. \\
\\
First of all, our gratitude sincerely goes to family who always give us dreams to be the excellent students.Our second gratitude belongs to our ad-visor Prof Madya Dr. Jusoh Yacob who has patiently trained and taught us to be the real researcher. \\
\\
In addition, we would like to thank to our lecturers that taught us and all of statistics student who have contribute and sharing the knowledge with us. We realize truly that this final year project needs the constructive criticism to be remarkable research findings. }
\pagebreak



%{\clearpage\SingleSpacing
\setstretch{1}
\tableofcontents\clearpage
\listoftables\clearpage
\listoffigures\clearpage
\chapter{\textrm{Abstract}}
\setstretch{1}

In this study, an attempt has been made to study the factors affecting underpaid salary among fresh graduates of Bsc (Hons.) Statistics from UiTM Kota Bahru, Kelantan. There are five factors that contribute in this study which are demographic information, job mismatch, expectation salary, monthly expenses and qualifications (skills, academic and experience). Primary data set has been used in which 200 fresh graduates from alumni degree in Science Statistics UiTM Kota Bharu have been targeted. There are several methods used to conduct this study. As for the research design, case study method is the most suitable method to be used. Besides, self-administered questionnaire and mail method of survey is used to collect data from the respondents selected. The respondents were selected by using simple random sampling method after the sample size has been identified by using Roasoft software. Moreover, to analyze the data obtained a series of multiple logistic regression, correlation testing and descriptive analysis were performed. In addition, checking for multicollinearity and checking the goodness of fit for the model were also carried out before obtaining the final model. This study has measured the underpaid salary by two approaches which are monthly expenses and qualification. Fresh graduates that suffer from underpaid salary do not have enough salary to spend for the things they want. Moreover, based on the data collected, most organizations do not give salary based on experience and educational level. Our analysis shows the influence factors towards underpaid salary among fresh graduates. Last but not least, it is recommended to expand the study area and population besides, the study should be carry out in other universities.
\\


%\listofacronyms\clearpage
%\listofappendices\clearpage}
%\setlinespacing{1.5}%utk spacing 1.5
\renewcommand{\baselinestretch}{2.0} % spacing 1.5 next
\mainmatter
\include{intro}
\include{probstatment}
\include{objective}
\include{significant}
\include{scope}
\include{definition}
\include{literaturereview}
\include{methodology}
\include{keputusan}
\include{conclusion}
\include{benefitnscheduling}


\backmatter
%\bibliographystyle{apalike}
\setstretch{1}
\bibliography{ref}

%% Comment out this line if you do not have a separate glossary.
\printglossary
\include{appendix}
\end{document}
